\begin{myexercise} 
    Erstellen Sie eine Hierarchie von Klassen für geometrische Formen und Körper.
    \begin{itemize}
    \item [a)] Wir beginnen im Zweidimensionalen.
    \begin{enumerate}
        \item Erstellen Sie zunächst eine Basisklasse ZweiDForm. Implementieren Sie dann die Methoden berechne\_flaeche() und berechne\_umfang(). Diese Methoden sollen in der Basisklasse eine Fehlermeldung zurückgeben, dass die Methoden nicht implementiert sind.
        \item Erstellen Sie nun zwei verschiedene Unterklassen von ZweiDForm: Kreis und Rechteck. Überschreiben Sie für beide die Methoden berechne\_flaeche() und berechne\_umfang().
        \item Testen Sie Ihre Implementierung aus Schritt 1 und 2.
        \item Führen Sie nun die Unterklasse regelmaessiges\_n\_Eck (von ZweiDForm) hinzu. Überschreiben Sie für diese die Methoden berechne\_flaeche() und berechne\_umfang(). Hinweis: Die Fläche eines regelmäßigen n-Ecks ist $A=\frac{n\cdot a^2}{4\tan{\frac{180\circ}{n}}}$, wobei $a$ die Seitenlänge und $n$ die Anzahl der Ecken ist.
        \item Testen Sie auch diese Klasse.
    \end{enumerate}

    \item [b)] Wir wollen nun auch höherdimensionale Körper erlauben.
    \begin{enumerate}
        \item Erstellen Sie eine neue oberste Basisklasse namens GeometrischerKoerper. Der Konstruktor soll weiterhin den Namen der Form speichern und zusätzlich die Anzahl der Dimensionen des Körpers. Fügen Sie eine generische Methode info() hinzu, die den Namen und die Dimensionen des Körpers ausgibt.
        \item Passen Sie die Klasse ZweiDForm so an, dass sie eine Unterklasse von GeometrischerKoerper ist.
        \item Fügen Sie eine Klasse DreiDKoerper hinzu, die die Methoden berechne\_oberflaeche() und berechne\_volumen() hat. Diese Methoden sollen wieder in der Basisklasse eine Fehlermeldung zurückgeben, dass die Methoden nicht implementiert sind.
        \item Erstellen Sie mindestens zwei Klassen von 3-dimensionalen Körpern (z.B. Wuerfel, Kugel) und implementieren Sie die Methoden berechne\_oberflaeche() und berechne\_volumen().
    \end{enumerate}

    \item [c)]Erweitern Sie die Klassenhierarchie um \href{h[https://de.wikipedia.org/wiki/Platonischer_K%C3%B6rper](https://de.wikipedia.org/wiki/Platonischer_K%C3%B6rper)}{platonische Körper} und demonstrieren Sie, wie die berechne\_flaeche()-Methoden der 2D-Klassen für die Berechnung der Oberfläche der 3D-Körper wiederverwendet werden können.
    \begin{enumerate}
        \item Fügen Sie die Klasse PlatonischerKoerper als Unterklasse von DreiDKoerper hinzu. Diese Klasse sollte zusätzlich die Attribute grundflaeche und anzahl\_flaechen haben. Die Grundfläche ist bei einem platonischen Körper immer ein regelmäßiges n-Eck.
        \item Überschreiben Sie die Methode berechne\_oberflaeche(). Nutzen Sie die bereits vorhandene Methode berechne\_flaeche() des regelmäßigen n-Ecks.
        \item Testen Sie Ihre Implementierung.
    \end{enumerate}
\end{itemize}
\end{myexercise}